\documentclass[11pt,a4paper,openany]{ctexbook}

\usepackage[margin=2.2cm]{geometry}
\usepackage{enumitem}
\usepackage{booktabs}
\usepackage{xcolor}
\usepackage{hyperref}
\usepackage{tcolorbox}
\usepackage{listings}
\usepackage{fancyhdr}

\definecolor{codebg}{RGB}{248,248,248}
\definecolor{notebg}{RGB}{240,248,255}
\definecolor{warnbg}{RGB}{255,248,240}

\lstset{
  basicstyle=\ttfamily\small,
  backgroundcolor=\color{codebg},
  frame=single,
  breaklines=true,
  columns=fullflexible,
  keepspaces=true,
  showstringspaces=false,
  extendedchars=false
}
\lstdefinelanguage{yaml}{}

\newtcolorbox{note}[1][]{
  colback=notebg, colframe=blue!40!black, title=提示, #1
}
\newtcolorbox{warn}[1][]{
  colback=warnbg, colframe=orange!70!black, title=注意, #1
}

\pagestyle{fancy}
\fancyhf{}
\fancyhead[LE,RO]{\thepage}
\fancyhead[LO]{\nouppercase{\rightmark}}
\fancyhead[RE]{\nouppercase{\leftmark}}
\renewcommand{\headrulewidth}{0.4pt}

\title{Elasticsearch 迁移记录\\\large 7.17.7 $\rightarrow$ 8.19.23}
\author{}
\date{}

\begin{document}
\frontmatter
\maketitle

\begin{note}
本文记录从 Elasticsearch 7.17.7 集群迁移至 8.19.23 集群的完整过程。主体采用 Reindex 远程拉取方式，并在第 6 章汇总快照恢复、Logstash、elasticsearch-dump 等备选方案，供不同场景参考。
\end{note}

\tableofcontents
\mainmatter

\chapter{环境与迁移方案}

\section{源集群与目标集群}

\begin{center}
\begin{tabular}{@{}lll@{}}
\toprule
项目 & 源集群 & 目标集群 \\
\midrule
版本 & 7.17.7 & 8.19.23 \\
部署方式 & RPM & deb \\
地址 & 192.168.191.146:9200 & 192.168.191.212:9200 \\
协议 & HTTP & HTTPS + 认证 \\
节点数 & 1 & 1 \\
集群状态 & yellow（副本未分配） & green \\
\bottomrule
\end{tabular}
\end{center}

\section{本次采用的方案：Reindex 远程拉取}

版本跨度 7.17 到 8.19 在快照恢复的兼容范围内，且源集群为单节点、数据量极小（约 8.4 MB）。综合评估后选择 \textbf{Reindex 远程拉取}：

\begin{itemize}[leftmargin=1.4em]
  \item 不依赖共享存储，无需在源端创建快照仓库；
  \item 目标端可直接读取源端 HTTP 接口，配置简单；
  \item 数据量小，同步执行即可完成，无需异步任务监控；
  \item 避免将源集群潜在的旧索引问题带入目标集群。
\end{itemize}

\begin{note}
Reindex 是逻辑迁移，不会复制源索引的 Settings 和 Mapping。目标索引需提前创建，或依赖动态映射自动推断。其他方案的对比见第 6 章。
\end{note}

\chapter{迁移前检查}

\section{源集群版本与健康状态}

在源服务器上执行：

\begin{lstlisting}[language=bash]
curl -u elastic:'<源密码>' "http://192.168.191.146:9200"
curl -u elastic:'<源密码>' "http://192.168.191.146:9200/_cluster/health?pretty"
\end{lstlisting}

返回 \texttt{"number": "7.17.7"} 确认版本正确。源集群 \texttt{status} 为 \texttt{yellow}，单节点下副本分片未分配属正常，主分片全部可用，不影响 Reindex。

\section{索引清单与数据量}

\begin{lstlisting}[language=bash]
curl -u elastic:'<源密码>' \
  "http://192.168.191.146:9200/_cat/indices?v&h=index,health,status,docs.count,store.size,pri.store.size&s=store.size:desc"
\end{lstlisting}

本次迁移涉及 9 个业务索引，总数据量约 8.4 MB：

\begin{center}
\begin{tabular}{@{}lrr@{}}
\toprule
索引 & 文档数 & 主分片大小 \\
\midrule
knowledge\_index & 61 & 8.3 MB \\
fwt01 & 6 & 56.3 KB \\
service & 1 & 20.3 KB \\
api & 2 & 9.8 KB \\
apisix & 1 & 6.5 KB \\
v1 & 1 & 5.3 KB \\
v2 & 1 & 5.3 KB \\
v3 & 1 & 5.3 KB \\
auth & 1 & 3.9 KB \\
\midrule
合计 & 75 & 约 8.4 MB \\
\bottomrule
\end{tabular}
\end{center}

另有 \texttt{index.php} 与 \texttt{casa} 两个空索引（0 文档、Mapping 为空），迁移中跳过。

\section{旧索引排查}

7.x 到 8.x 迁移要求不存在 6.x 或更早版本创建的索引。检查创建时间：

\begin{lstlisting}[language=bash]
curl -u elastic:'<源密码>' \
  "http://192.168.191.146:9200/_cat/indices?v&h=index,creation.date.string&s=creation.date.string:asc"
\end{lstlisting}

本次最早的索引创建于 2023-01-12，全部晚于 7.0 发布时间（2019 年 4 月），无旧索引需要预处理。

\section{源集群监听与连通性}

确认源集群 \texttt{network.host} 绑定地址：

\begin{lstlisting}[language=bash]
grep -E "network.host|http.port" /etc/elasticsearch/elasticsearch.yml
\end{lstlisting}

本次源集群配置为 \texttt{network.host: 0.0.0.0}，监听所有网卡。在目标机器上测试连通性：

\begin{lstlisting}[language=bash]
curl -u elastic:'<源密码>' "http://192.168.191.146:9200"
\end{lstlisting}

能返回 7.17.7 版本信息即表示网络与认证均通。

\chapter{目标集群配置}

\section{配置远程白名单}

编辑目标集群 \texttt{/etc/elasticsearch/elasticsearch.yml}，添加源集群地址。注意只写 \texttt{host:port}，不带协议前缀：

\begin{lstlisting}[language=yaml]
reindex.remote.whitelist: "192.168.191.146:9200"
\end{lstlisting}

保存后重启目标集群：

\begin{lstlisting}[language=bash]
sudo systemctl restart elasticsearch
\end{lstlisting}

重启后验证服务正常：

\begin{lstlisting}[language=bash]
curl --cacert /etc/elasticsearch/certs/http_ca.crt \
  -u elastic:'<目标密码>' https://localhost:9200
\end{lstlisting}

\section{源集群为 HTTPS 时的额外配置}

本次源集群为 HTTP，无需额外配置。若源集群启用了 HTTPS，需在目标集群补充 CA 证书路径：

\begin{lstlisting}[language=yaml]
reindex.ssl.certificate_authorities: "/path/to/源集群-ca.crt"
\end{lstlisting}

仅在测试环境可临时关闭证书校验：

\begin{lstlisting}[language=yaml]
reindex.ssl.verification_mode: none
\end{lstlisting}

\chapter{索引 Mapping 与预建}

\section{获取源索引 Mapping}

\begin{lstlisting}[language=bash]
curl -u elastic:'<源密码>' \
  "http://192.168.191.146:9200/knowledge_index/_mapping?pretty"
\end{lstlisting}

7.x 返回结构外层包含索引名与 \texttt{mappings}，创建索引时只需 \texttt{mappings} 内部内容，不需要外层索引名包裹。

\section{需要预建的索引}

含 \texttt{date}、\texttt{long} 等显式类型的索引必须预建，否则动态映射可能将日期推断为 text、数值推断为 keyword，导致查询行为不一致。

\begin{center}
\begin{tabular}{@{}lll@{}}
\toprule
索引 & 需预建原因 & 关键字段 \\
\midrule
knowledge\_index & date 字段 & createTime, updateTime \\
fwt01 & date + long & create\_time, update\_time, deleted \\
service & long 字段 & data.limit, data.page, data.start, tid \\
apisix & long 字段 & timeout \\
\bottomrule
\end{tabular}
\end{center}

其余索引（\texttt{api}、\texttt{v1}、\texttt{v2}、\texttt{v3}、\texttt{auth}）字段仅含 text + keyword 子字段与 object，可由动态映射正确推断，无需预建。

\section{预建索引示例}

以 \texttt{knowledge\_index} 为例，单节点部署下副本数设为 0，避免未分配副本导致索引保持 yellow：

\begin{lstlisting}[language=bash]
curl --cacert /etc/elasticsearch/certs/http_ca.crt \
  -u elastic:'<目标密码>' \
  -XPUT "https://localhost:9200/knowledge_index" \
  -H "Content-Type: application/json" \
  -d '{
  "settings": {"number_of_shards": 1, "number_of_replicas": 0},
  "mappings": {
    "properties": {
      "createTime": {"type": "date"},
      "updateTime": {"type": "date"},
      "content": {"type": "text",
        "fields": {"keyword": {"type": "keyword", "ignore_above": 256}}}
    }
  }
}'
\end{lstlisting}

\texttt{fwt01}、\texttt{service}、\texttt{apisix} 按各自 Mapping 同样预建，返回 \texttt{"acknowledged": true} 即成功。

\chapter{执行 Reindex}

\section{单索引 Reindex}

以 \texttt{knowledge\_index} 为例，数据量小，同步等待完成：

\begin{lstlisting}[language=bash]
curl --cacert /etc/elasticsearch/certs/http_ca.crt \
  -u elastic:'<目标密码>' \
  -XPOST "https://localhost:9200/_reindex?wait_for_completion=true" \
  -H "Content-Type: application/json" \
  -d '{
  "source": {
    "remote": {
      "host": "http://192.168.191.146:9200",
      "username": "elastic",
      "password": "<源密码>"
    },
    "index": "knowledge_index"
  },
  "dest": {"index": "knowledge_index"}
}'
\end{lstlisting}

成功返回示例：

\begin{lstlisting}
{"took":1277,"timed_out":false,"total":61,"created":61,
 "failures":[]}
\end{lstlisting}

\section{多索引批量 Reindex}

\texttt{\_reindex} 的 \texttt{dest.index} 不支持数组，一次请求只能写入一个目标索引。批量迁移应使用循环逐个执行：

\begin{lstlisting}[language=bash]
for idx in fwt01 service api apisix v1 v2 v3 auth; do
  echo "=== Reindexing $idx ==="
  curl --cacert /etc/elasticsearch/certs/http_ca.crt \
    -u elastic:'<目标密码>' \
    -XPOST "https://localhost:9200/_reindex?wait_for_completion=true" \
    -H "Content-Type: application/json" \
    -d "{
      \"source\": {
        \"remote\": {
          \"host\": \"http://192.168.191.146:9200\",
          \"username\": \"elastic\",
          \"password\": \"<源密码>\"
        },
        \"index\": \"$idx\"
      },
      \"dest\": {\"index\": \"$idx\"}
    }"
  echo
done
\end{lstlisting}

\begin{warn}
若将 \texttt{source.index} 写成数组而 \texttt{dest.index} 写固定占位名（如 \texttt{PLACEHOLDER}），所有文档都会因目标索引名非法而失败，返回 \texttt{invalid\_index\_name\_exception}。此失败不会造成数据损坏，但需改用循环方式重新执行。
\end{warn}

\chapter{其他迁移方案}

本章汇总 Elasticsearch 常用的其他迁移方式，说明各自的版本限制、适用场景与基本操作，便于在不同环境中选型。

\section{方案对比}

\begin{center}
\begin{tabular}{@{}p{0.18\textwidth}p{0.28\textwidth}p{0.25\textwidth}p{0.2\textwidth}@{}}
\toprule
方案 & 适用场景 & 版本限制 & 对源端影响 \\
\midrule
快照恢复 & 同版本或相邻主版本，大数据量 & 不能跨一个以上主版本；目标版本须 $\ge$ 源版本[citation:3] & 低 \\
Logstash & 跨多个主版本，需要数据转换或过滤 & 无严格限制，Logstash 主版本建议与目标集群一致[citation:6] & 中高 \\
elasticsearch-dump & 小数据量，按需选择 settings/mapping/data & 无严格限制 & 中 \\
Reindex 远程拉取 & 跨版本，源端 HTTP 可达，无需共享存储 & 无严格限制 & 中 \\
\bottomrule
\end{tabular}
\end{center}

\section{快照与恢复}

快照是最快、最可靠的迁移方式，但版本约束严格：只能恢复到相同或更高版本的集群，且不能跨一个以上主版本。例如 6.x 快照可以恢复到 7.x，但不能直接恢复到 8.x[citation:3]。7.17 到 8.19 在允许范围内，但需要源端和目标端能访问同一个快照仓库（如 S3、NFS）。

基本流程：

\begin{lstlisting}[language=bash]
# 在源集群创建快照
PUT _snapshot/my_repo/snapshot_1?wait_for_completion=true

# 在目标集群注册同一仓库并恢复
POST _snapshot/my_repo/snapshot_1/_restore
\end{lstlisting}

\begin{warn}
目标集群注册仓库时建议设为只读（\texttt{readonly: true}），防止 8.x 集群意外写入导致快照格式混乱。若源集群 Upgrade Assistant 提示有需 Reindex 的旧索引，快照恢复后这些索引在 8.x 中依然不可用。
\end{warn}

\section{Logstash}

Logstash 支持从源集群读取数据并写入目标集群，适合跨多个主版本或需要对数据做过滤、转换的场景[citation:6]。需要注意：

\begin{itemize}[leftmargin=1.4em]
  \item Logstash 主版本建议与目标 Elasticsearch 集群一致[citation:6]；
  \item 跨大版本时需注意索引 type 差异，7.x 以上已移除多 type 支持；
  \item 部署 Logstash 的机器需要能同时访问源集群和目标集群。
\end{itemize}

典型的 Logstash 配置（\texttt{es-to-es.conf}）结构：

\begin{lstlisting}
input {
  elasticsearch {
    hosts => ["http://192.168.191.146:9200"]
    user => "elastic"
    password => "<源密码>"
    index => "knowledge_index"
    query => '{"query": {"match_all": {}}}'
  }
}
output {
  elasticsearch {
    hosts => ["https://192.168.191.212:9200"]
    user => "elastic"
    password => "<目标密码>"
    ssl => true
    cacert => "/etc/elasticsearch/certs/http_ca.crt"
    index => "knowledge_index"
  }
}
\end{lstlisting}

\section{elasticsearch-dump}

elasticsearch-dump 是基于 Node.js 的开源工具，适合小数据量迁移，可以按类型分别迁移 settings、mapping、data[citation:6]。安装：

\begin{lstlisting}[language=bash]
npm install elasticdump -g
\end{lstlisting}

分步迁移单个索引：

\begin{lstlisting}[language=bash]
# 先迁 settings 和 mapping，再迁数据
elasticdump --input=http://192.168.191.146:9200/knowledge_index \
  --output=http://192.168.191.212:9200/knowledge_index --type=settings
elasticdump --input=http://192.168.191.146:9200/knowledge_index \
  --output=http://192.168.191.212:9200/knowledge_index --type=mapping
elasticdump --input=http://192.168.191.146:9200/knowledge_index \
  --output=http://192.168.191.212:9200/knowledge_index --type=data
\end{lstlisting}

若目标集群启用 HTTPS，需加 \texttt{-ca} 参数指定证书；若密码含特殊字符，需先做 URL 编码[citation:6]。

\section{选型建议}

本次迁移数据量极小（8.4 MB）、源端 HTTP 可达，Reindex 远程拉取是最直接的选择。若未来数据量增大到 GB 级以上且版本在快照兼容范围内，可优先考虑快照恢复以缩短迁移窗口。若需要跨多个主版本或对字段做清洗转换，Logstash 更合适。elasticsearch-dump 适合临时、小批量的按需迁移，不适合作为大规模生产迁移的主力方案。

\chapter{迁移结果验证}

\section{文档数对比}

\begin{lstlisting}[language=bash]
curl --cacert /etc/elasticsearch/certs/http_ca.crt \
  -u elastic:'<目标密码>' \
  "https://localhost:9200/_cat/indices?v&h=index,docs.count,store.size&s=index:asc"
\end{lstlisting}

本次迁移结果与源端逐项一致：

\begin{center}
\begin{tabular}{@{}lrrl@{}}
\toprule
索引 & 源端文档数 & 目标端文档数 & 状态 \\
\midrule
knowledge\_index & 61 & 61 & 一致 \\
fwt01 & 6 & 6 & 一致 \\
service & 1 & 1 & 一致 \\
api & 2 & 2 & 一致 \\
apisix & 1 & 1 & 一致 \\
v1 & 1 & 1 & 一致 \\
v2 & 1 & 1 & 一致 \\
v3 & 1 & 1 & 一致 \\
auth & 1 & 1 & 一致 \\
\midrule
合计 & 75 & 75 & 全部一致 \\
\bottomrule
\end{tabular}
\end{center}

\section{集群健康状态}

\begin{lstlisting}[language=bash]
curl --cacert /etc/elasticsearch/certs/http_ca.crt \
  -u elastic:'<目标密码>' \
  "https://localhost:9200/_cluster/health?pretty"
\end{lstlisting}

若 \texttt{knowledge\_index} 副本数设为 1 而目标为单节点，索引会保持 yellow。可将副本数调为 0 使集群显示 green：

\begin{lstlisting}[language=bash]
curl --cacert /etc/elasticsearch/certs/http_ca.crt \
  -u elastic:'<目标密码>' \
  -XPUT "https://localhost:9200/knowledge_index/_settings" \
  -H "Content-Type: application/json" \
  -d '{"index": {"number_of_replicas": 0}}'
\end{lstlisting}

\section{数据抽查}

随机取一条文档确认字段内容完整：

\begin{lstlisting}[language=bash]
curl --cacert /etc/elasticsearch/certs/http_ca.crt \
  -u elastic:'<目标密码>' \
  "https://localhost:9200/knowledge_index/_search?size=1&pretty"
\end{lstlisting}

对 date 字段做范围查询，确认 Mapping 预建生效：

\begin{lstlisting}[language=bash]
curl --cacert /etc/elasticsearch/certs/http_ca.crt \
  -u elastic:'<目标密码>' \
  -XGET "https://localhost:9200/knowledge_index/_count" \
  -H "Content-Type: application/json" \
  -d '{"query": {"range": {"createTime": {"gte": "2020-01-01"}}}}'
\end{lstlisting}

\chapter{收尾工作}

\section{移除远程白名单}

迁移完成后，若目标集群不再需要从源集群拉数据，建议移除白名单并重启，减少暴露面：

\begin{lstlisting}[language=yaml]
# reindex.remote.whitelist: "192.168.191.146:9200"
\end{lstlisting}

\section{增量数据同步}

Reindex 是快照式拉取。若源集群在迁移期间仍有写入，业务切换前需补充同步。两种方式：

\begin{itemize}[leftmargin=1.4em]
  \item 基于时间范围对每个索引做一次增量 \texttt{\_reindex}，使用 \texttt{query} 限定新增文档；
  \item 短暂停写后重跑一次全量 Reindex，再切换业务。
\end{itemize}

增量 Reindex 示例：

\begin{lstlisting}[language=bash]
curl --cacert /etc/elasticsearch/certs/http_ca.crt \
  -u elastic:'<目标密码>' \
  -XPOST "https://localhost:9200/_reindex?wait_for_completion=true" \
  -H "Content-Type: application/json" \
  -d '{
  "source": {
    "remote": {
      "host": "http://192.168.191.146:9200",
      "username": "elastic",
      "password": "<源密码>"
    },
    "index": "knowledge_index",
    "query": {"range": {"updateTime": {"gte": "2026-10-28T00:00:00Z"}}}
  },
  "dest": {"index": "knowledge_index"}
}'
\end{lstlisting}

\section{应用连接切换}

业务端连接从源集群切换到目标集群时，需注意协议差异：

\begin{itemize}[leftmargin=1.4em]
  \item 源端为 HTTP，目标端为 HTTPS + 认证；
  \item 客户端需配置 CA 证书（\texttt{http\_ca.crt}）；
  \item 建议创建专用用户或 API Key，而非直接使用 \texttt{elastic} 超级用户。
\end{itemize}

\chapter{迁移要点总结}

\begin{enumerate}[leftmargin=1.4em]
  \item 迁移前先确认源集群版本、索引清单、创建时间，排除 6.x 旧索引。
  \item 确认源集群监听地址与目标端连通性，避免白名单配置后才发现网络不通。
  \item Reindex 不复制 Mapping，含 date/long 等显式类型的索引必须预建。
  \item 单节点部署下副本数设为 0，避免索引长期处于 yellow。
  \item \texttt{\_reindex} 的 \texttt{dest.index} 不支持数组，批量迁移用循环逐个执行。
  \item 其他方案各有约束：快照有版本跨度限制，Logstash 需注意 type 和版本匹配，elasticsearch-dump 适合小批量按需迁移。
  \item 迁移后逐项对比文档数，并对 date 字段做范围查询验证 Mapping 生效。
  \item 收尾时移除白名单、处理增量数据、切换应用连接并改用专用凭据。
\end{enumerate}

\end{document}